\section{Externí paměti a~úložiště}\label{sec:hw-externi-pameti}

Hlavní paměť uchovává data pouze po dobu, kdy je počítač napájen. Programy, dokumenty a~další soubory proto potřebujeme ukládat do nevolatilní paměti, jejíž obsah zůstane zachován i~po vypnutí. Mezi nejběžnější dlouhodobá úložiště patří pevné disky HDD a~disky SSD. Obě zařízení poskytují operačnímu systému adresovatelný prostor rozdělený do bloků, ale fyzikální princip jejich fungování je odlišný.

\begin{figure}[ht]
    \centering
    \begin{tikzpicture}[
    x=1cm,
    y=1cm,
    >=Stealth,
    every node/.style={font=\small},
    panel/.style={
        draw,
        thick,
        rounded corners,
        minimum width=6.0cm,
        minimum height=3.8cm,
        align=center
    }
]
    \node[panel, fill=blue!8] (hdd) at (-3.6,0) {};
    \node[font=\large\bfseries] at (-3.6,1.35) {HDD};
    \draw[thick, fill=gray!18] (-4.2,0.25) circle (0.85);
    \draw[thick, fill=white] (-4.2,0.25) circle (0.17);
    \draw[thick] (-4.2,0.25) circle (0.52);
    \draw[thick, fill=gray!40] (-2.55,0.85) circle (0.14);
    \draw[very thick] (-2.55,0.85) -- (-3.85,0.45);
    \draw[thick, fill=black] (-3.85,0.45) circle (0.06);
    \node[align=center, font=\scriptsize] at (-3.6,-1.15)
        {magnetické plotny\\ mechanický pohyb};

    \node[panel, fill=green!9] (ssd) at (3.6,0) {};
    \node[font=\large\bfseries] at (3.6,1.35) {SSD};
    \draw[thick, rounded corners=3pt, fill=green!35]
        (1.9,-0.25) rectangle (5.1,0.65);
    \foreach \x in {2.45,3.25,4.05}{
        \draw[thick, rounded corners=2pt, fill=black!72]
            (\x-0.30,-0.02) rectangle (\x+0.30,0.40);
    }
    \draw[thick, rounded corners=2pt, fill=blue!40]
        (4.55,-0.02) rectangle (4.95,0.40);
    \draw[thick, fill=yellow!65]
        (5.1,-0.10) rectangle (5.42,0.50);
    \node[align=center, font=\scriptsize] at (3.6,-1.15)
        {flash paměť\\ bez pohyblivých částí};

    \node[
        draw,
        thick,
        rounded corners,
        fill=yellow!16,
        minimum width=9.6cm,
        minimum height=0.75cm,
        align=center
    ] at (0,-2.75)
        {Obě zařízení uchovávají data i~po vypnutí napájení.};
\end{tikzpicture}

    \caption{Základní konstrukční rozdíl mezi HDD a~SSD.}
    \label{fig:hw-hdd-ssd}
\end{figure}

HDD ukládá data magneticky na rotující plotny a~používá pohyblivé čtecí a~zapisovací hlavy. SSD ukládá data elektronicky do buněk flash paměti a~neobsahuje části, které by se při běžném čtení mechanicky pohybovaly. Díky tomu má SSD zpravidla podstatně rychlejší náhodný přístup, je odolnější vůči otřesům a~pracuje tiše. HDD naopak často nabízí velkou kapacitu za nižší cenu.

\subsection{Pevný disk HDD}\label{subsec:hw-hdd}

Uvnitř HDD se nachází jedna nebo více kruhových ploten upevněných na společném vřetenu. Povrch každé plotny je pokryt magnetickou vrstvou. Plotny se při provozu otáčejí a~nad jejich povrchem se pohybují čtecí a~zapisovací hlavy umístěné na společném rameni. Hlava se povrchu za běžného provozu nedotýká; změny magnetizace povrchu používá ke čtení a~zápisu bitů.

\begin{figure}[ht]
    \centering
    \begin{tikzpicture}[
    x=1cm,
    y=1cm,
    >=Stealth,
    every node/.style={font=\small},
    label/.style={font=\scriptsize, fill=white, inner sep=1.5pt},
    callout/.style={->, thick}
]
    \node[font=\bfseries] at (-3.8,2.95) {Pohled shora};
    \draw[thick, rounded corners=5pt, fill=blue!5]
        (-7.0,-2.4) rectangle (-0.6,2.55);
    \draw[thick, fill=gray!18] (-4.65,0.05) circle (1.75);
    \foreach \r in {0.55,0.95,1.35}{
        \draw[gray!55] (-4.65,0.05) circle (\r);
    }
    \draw[thick, fill=white] (-4.65,0.05) circle (0.30);
    \draw[thick, fill=gray!45] (-1.65,0.85) circle (0.24);
    \draw[very thick] (-1.65,0.85) -- (-3.85,0.32);
    \draw[thick, fill=black] (-3.85,0.32) rectangle (-3.65,0.46);

    \node[label, anchor=east] (platterlabel) at (-6.3,2.05) {plotna};
    \draw[callout] (platterlabel.east) -- (-5.65,1.42);
    \node[label] (armlabel) at (-2.0,2.05) {rameno};
    \draw[callout] (armlabel.south) -- (-2.45,0.67);
    \node[label] (headlabel) at (-2.9,-1.85) {čtecí/zapisovací hlava};
    \draw[callout] (headlabel.north) -- (-3.72,0.38);

    \draw[gray!45, thick] (0,-2.6) -- (0,2.8);

    \node[font=\bfseries] at (3.9,2.95) {Pohled ze strany};
    \draw[thick, rounded corners=5pt, fill=blue!5]
        (0.7,-2.4) rectangle (7.1,2.55);

    \foreach \y in {-1.2,0,1.2}{
        \draw[thick, fill=gray!18] (3.3,\y) ellipse (2.0 and 0.33);
        \draw[gray!55] (3.3,\y) ellipse (1.1 and 0.18);
    }
    \draw[very thick, blue!65!black] (3.3,-1.75) -- (3.3,1.75);

    \draw[thick, fill=gray!45] (6.35,-1.75) rectangle (6.7,1.75);
    \foreach \y in {-1.2,0,1.2}{
        \draw[very thick] (6.35,\y+0.18) -- (4.25,\y+0.05);
        \draw[thick, fill=black] (4.15,\y-0.02) rectangle (4.35,\y+0.12);
    }

    \node[label, anchor=west] (spindlelabel) at (0.95,2.0) {vřeteno};
    \draw[callout] (spindlelabel.east) -- (3.15,1.55);
    \node[label, anchor=west] (stacklabel) at (0.95,-2.0) {soubor ploten};
    \draw[callout] (stacklabel.east) -- (2.2,-1.2);
    \node[label] (assemblylabel) at (5.7,2.05) {mechanismus ramen};
    \draw[callout] (assemblylabel.south) -- (6.5,1.4);
\end{tikzpicture}

    \caption{Zjednodušený pohled na vnitřní strukturu HDD shora a~ze strany.}
    \label{fig:hw-hdd-struktura}
\end{figure}

Data jsou na povrchu plotny uspořádána do soustředných \emph{stop}. Každá stopa se dále dělí na \emph{sektory}. Stopy ve stejné vzdálenosti od středu na různých površích ploten tvoří \emph{cylindr}. V~moderním disku se operační systém nemusí zabývat přesnou fyzickou polohou dat; řadič disku převádí logická čísla bloků na konkrétní místa uvnitř zařízení. Geometrický model je však stále užitečný k~pochopení výkonu HDD.

\begin{figure}[ht]
    \centering
    \begin{tikzpicture}[
    scale=0.95,
    transform shape,
    x=1cm,
    y=1cm,
    >=Stealth,
    every node/.style={font=\small},
    track/.style={draw, gray!60},
    sectorline/.style={draw, thin, gray!50},
    callout/.style={->, thick},
    label/.style={font=\small, fill=white, inner sep=2pt}
]
    \draw[thick, fill=gray!10] (0,0) circle (3.0);
    \path[draw=orange!80!black, thick, fill=yellow!55]
        (25:1.55) arc (25:55:1.55) --
        (55:2.1) arc (55:25:2.1) -- cycle;

    \foreach \radius in {0.8,1.15,1.55,1.85,2.1,2.4,2.7}{
        \draw[track] (0,0) circle (\radius);
    }
    \foreach \angle in {0,30,...,330}{
        \draw[sectorline] (\angle:0.62) -- (\angle:3.0);
    }
    \draw[thick, fill=white] (0,0) circle (0.62);
    \draw[thick, fill=white] (0,0) circle (0.20);

    \node[label, anchor=east] (platter) at (-4.0,2.45) {\textbf{plotna}};
    \draw[callout] (platter.east) -- (-2.55,1.55);

    \node[label, anchor=east] (tracklabel) at (-4.0,1.05) {\textbf{stopa}};
    \draw[callout] (tracklabel.east) -- (-1.65,0.8);

    \node[label, anchor=east] (spindle) at (-4.0,-0.35) {\textbf{vřeteno}};
    \draw[callout] (spindle.east) -- (-0.45,-0.15);

    \node[label, anchor=west] (sector) at (4.0,1.9) {\textbf{sektor}};
    \draw[callout] (sector.west) -- (1.55,1.3);

    \draw[->, very thick, blue!70!black]
        (2.55,2.0) arc (38:-25:3.25);
    \node[font=\scriptsize, text=blue!70!black] at (4.05,-1.35) {směr otáčení};

    \node[
        draw,
        thick,
        rounded corners,
        fill=blue!7,
        align=center,
        text width=11.0cm
    ] at (0,-4.0)
        {Povrch plotny je rozdělen na soustředné \textbf{stopy}; každá stopa se dále dělí na \textbf{sektory}.};
\end{tikzpicture}

    \caption{Stopy a~sektory na jedné plotně HDD.}
    \label{fig:hw-hdd-plotna}
\end{figure}

Při čtení bloku musí disk nejprve přesunout rameno nad správnou stopu. Doba tohoto pohybu se označuje jako \emph{doba vystavení} neboli \emph{seek time}. Poté disk čeká, než se požadovaný sektor otočí pod hlavu; vzniká \emph{rotační zpoždění}. Teprve potom lze data skutečně přenést. Přesun hlavy a~čekání na natočení plotny jsou mechanické operace, a~proto jsou ve srovnání s~elektronickými operacemi pomalé.

Sekvenční čtení velkého souboru je pro HDD výhodné, protože sousední sektory procházejí pod hlavou postupně a~rameno se nemusí často přesouvat. Naopak mnoho malých požadavků na vzdálená místa disku vede k~opakovanému pohybu hlav a~výkon výrazně klesá. Rozdíl mezi sekvenčním a~náhodným přístupem je proto u~HDD zvlášť výrazný.

\warningbox{Při prudkém nárazu může dojít k~poškození mechanických částí HDD nebo povrchu plotny. Běžná záloha proto nemá být uložena pouze na druhém oddílu téhož fyzického disku; porucha zařízení by mohla zasáhnout původní data i~jejich „zálohu“.}

\subsection{Disk SSD}\label{subsec:hw-ssd}

SSD neukládá data magneticky, ale do buněk nevolatilní flash paměti. Kromě paměťových čipů obsahuje řadič, který přijímá požadavky počítače, vyhledává fyzické umístění dat, opravuje některé chyby a~rozděluje zápisy mezi paměťové buňky. Operační systém přitom stejně jako u~HDD pracuje s~logickými bloky a~nemusí znát konkrétní uspořádání buněk uvnitř zařízení.

\begin{figure}[ht]
    \centering
    \begin{tikzpicture}[
    x=1cm,
    y=1cm,
    >=Stealth,
    every node/.style={font=\small},
    block/.style={
        draw,
        thick,
        rounded corners,
        align=center
    },
    page/.style={
        draw,
        thick,
        minimum width=1.05cm,
        minimum height=0.72cm,
        align=center,
        font=\scriptsize
    },
    arrow/.style={->, thick}
]
    \node[
        block,
        fill=blue!10,
        minimum width=2.7cm,
        minimum height=1.0cm
    ] (computer) at (-6.0,0.8) {Počítač\\ \scriptsize logické bloky};

    \node[
        block,
        fill=orange!18,
        minimum width=3.0cm,
        minimum height=2.4cm
    ] (controller) at (-2.3,0.8) {\textbf{Řadič SSD}\\[1mm]\scriptsize převod adres\\ oprava chyb\\ vyrovnávání opotřebení};

    \node[
        block,
        fill=green!8,
        minimum width=6.2cm,
        minimum height=4.4cm
    ] (flash) at (4.0,0.8) {};
    \node[font=\bfseries] at (4.0,2.45) {Flash paměť};

    \node[
        draw,
        thick,
        rounded corners,
        fill=green!16,
        minimum width=5.1cm,
        minimum height=2.4cm
    ] at (4.0,0.55) {};
    \node[font=\small\bfseries] at (4.0,1.45) {mazací blok};

    \node[page, fill=white] at (2.15,0.35) {stránka 0};
    \node[page, fill=white] at (3.4,0.35) {stránka 1};
    \node[page, fill=yellow!50] at (4.65,0.35) {stránka 2};
    \node[page, fill=white] at (5.9,0.35) {stránka 3};

    \node[
        block,
        fill=gray!10,
        minimum width=5.1cm,
        minimum height=0.65cm,
        font=\scriptsize
    ] at (4.0,-0.85) {další bloky a~náhradní prostor};

    \draw[arrow] (computer.east) -- (controller.west);
    \draw[arrow] (controller.east) -- node[above, align=center] {\scriptsize fyzické\\ adresy} (flash.west);
\end{tikzpicture}

    \caption{Zjednodušené části SSD a~organizace flash paměti.}
    \label{fig:hw-ssd-struktura}
\end{figure}

Flash paměť čte a~zapisuje data po \emph{stránkách}, ale mazání probíhá po větších \emph{blocích} obsahujících více stránek. Již zapsanou stránku proto nelze libovolně přepsat stejným způsobem jako bajt v~RAM. Řadič obvykle zapíše novou verzi dat na jiné volné místo, změní svou převodní tabulku a~později uvolní bloky obsahující již neplatná data.

Každá paměťová buňka snese pouze omezený počet mazacích cyklů. Řadič proto používá \emph{vyrovnávání opotřebení}, kterým se snaží rozložit zápisy mezi různé fyzické bloky. Současně může mít SSD část kapacity vyhrazenu jako náhradní prostor. Tyto mechanismy prodlužují životnost zařízení a~umožňují skrýt složitější vlastnosti flash paměti za jednoduché rozhraní logických bloků.

\notebox{SSD nemá mechanické vyhledávání stopy ani rotační zpoždění, a~proto je jeho náhodný přístup podstatně rychlejší než u~HDD. Přesto ani SSD není stejně rychlé jako RAM a~při zápisu musí jeho řadič řešit omezení flash paměti.}

\subsection{Porovnání HDD a~SSD}

Volba úložiště závisí na požadované kapacitě, rychlosti, ceně i~způsobu použití. Pro systémový disk a~často spouštěné aplikace je obvykle výhodné SSD, protože zkracuje čekání při mnoha menších náhodných přístupech. HDD se může hodit pro velké archivy nebo zálohy, u~nichž je důležitá kapacita a~převládá sekvenční přenos.

\begin{table}[ht]
    \centering
    \begin{tabular}{|l|c|c|}
        \hline
        \textbf{Vlastnost} & \textbf{HDD} & \textbf{SSD} \\ \hline
        Princip uložení & magnetický & elektronický \\ \hline
        Pohyblivé části & ano & ne \\ \hline
        Náhodný přístup & pomalejší & rychlejší \\ \hline
        Hlučnost & slyšitelný pohyb & tiché \\ \hline
        Typická přednost & kapacita a~cena & rychlost a~odolnost \\ \hline
        Uchování bez napájení & ano & ano \\ \hline
    \end{tabular}
    \caption{Obecné porovnání vlastností HDD a~SSD.}
    \label{tab:hw-hdd-ssd}
\end{table}

Tabulka~\ref{tab:hw-hdd-ssd} zachycuje obecné tendence, nikoliv záruku pro každý konkrétní model. Rychlé HDD může v~některé úloze překonat velmi pomalé SSD a~výsledný výkon ovlivňuje také použité rozhraní, řadič, velikost požadavků nebo zaplnění zařízení. Při praktickém porovnávání je proto potřeba sledovat parametry a~měření konkrétních výrobků.
